Zeichnen Sie das Zustandsdiagramm einer Turing-Maschine, die eine
positive Binärzahl dekrementiert.
\begin{teilaufgaben}
\item
Ermitteln Sie die Berechnungsgeschichte für die Verarbeitung des Inputwortes
\texttt{1100}.
\item
Ermitteln Sie die Berechnungsgeschichte für die Verarbeitung des Inputwortes
\texttt{1}.
\item
Ermitteln Sie die Berechnungsgeschichte für die Verarbeitung des Inputwortes
\texttt{0}.
\end{teilaufgaben}

\begin{hinweis}
In der Aufgabe wird nur verlangt, dass die Turing-Maschine {\em positive}
Binärzahlen dekrementieren können muss.
In Teilaufgabe c) wird also nach dem Verhalten der Turing-Maschine in einem
Fall gefragt, in dem ein korrektes Resultat nicht erwartet wird.
\end{hinweis}

\begin{loesung}
Das folgende Zustandsdiagramm löst das Problem:
\begin{center}
\def\l{3}
\begin{tikzpicture}[>=latex,thick]
\coordinate (q0) at (0,0);
\coordinate (q1) at (\l,0);
\coordinate (qa) at ({2*\l},0);
\coordinate (qr) at (\l,-2);
\draw (q0) circle[radius=0.35cm];
\draw (q1) circle[radius=0.35cm];
\draw (qa) ellipse(0.65cm and 0.35cm);
\draw (qa) ellipse(0.6cm and 0.3cm);
\draw (qr) ellipse(0.65cm and 0.35cm);
\draw (qr) ellipse(0.6cm and 0.3cm);
\node at (q0) {$q_0$};
\node at (q1) {$q_1$};
\node at (qa) {$q_{\text{accept}}$};
\node at (qr) {$q_{\text{reject}}$};
\draw[->,shorten >= 0.35cm] ($(q0)+(-1.5,0)$) -- (q0);

\draw[->,shorten <= 0.35cm,shorten >= 0.35cm,distance=1cm]
	(q0) to[out=60,in=120,distance=1cm] (q0);
\node at ($(q0)+(0,1.2)$) {$\texttt{0}\to\texttt{0},R
\atop
\texttt{1}\to\texttt{1},R$};

\draw[->,shorten <= 0.35cm,shorten >= 0.35cm] (q0) -- (q1);
\node at ($0.5*(q0)+0.5*(q1)$) [above] {$\scriptstyle\blank\to\blank,L$};

\draw[->,shorten <= 0.35cm,shorten >= 0.35cm,distance=1cm]
	(q1) to[out=60,in=120,distance=1cm] (q1);
\node at ($(q1)+(0,1.1)$) {$\scriptstyle\texttt{0}\to\texttt{1},L$};

\draw[->,shorten <= 0.35cm,shorten >= 0.35cm] (q1) -- (qr);
\node at ($0.4*(q1)+0.5*(qr)$) [left] {$\scriptstyle\blank\to\blank,R$};

\draw[->,shorten <= 0.35cm,shorten >= 0.65cm] (q1) -- (qa);
\node at ($0.5*(q1)+0.5*(qa)$) [above]
	{$\scriptstyle\texttt{1}\to\texttt{0},R$};

\end{tikzpicture}
\end{center}
\def\z#1{\hbox to0.4cm{\hfill\clap{$#1$}\hfill}}
\def\zb{\hbox to0.4cm{\hfill\clap{\blank}\hfill}}
\def\zt#1{\hbox to0.4cm{\hfill\clap{\texttt{#1}}\hfill}}
\begin{teilaufgaben}
\item
Die Verarbeitung des Wortes $\texttt{1100}_2=12_{10}$ ergibt
\begin{align*}
&\zb \zb \z{q_0} \zt{1}  \zt{1}  \zt{0}  \zt{0}  \zb \zb \\
&\zb \zb \zt{1}  \z{q_0} \zt{1}  \zt{0}  \zt{0}  \zb \zb \\
&\zb \zb \zt{1}  \zt{1}  \z{q_0} \zt{0}  \zt{0}  \zb \zb \\
&\zb \zb \zt{1}  \zt{1}  \zt{0}  \z{q_0} \zt{0}  \zb \zb \\
&\zb \zb \zt{1}  \zt{1}  \zt{0}  \zt{0}  \z{q_0} \zb \zb \\
&\zb \zb \zt{1}  \zt{1}  \zt{0}  \z{q_1} \zt{0}  \zb \zb \\
&\zb \zb \zt{1}  \zt{1}  \z{q_1} \zt{0}  \zt{1}  \zb \zb \\
&\zb \zb \zt{1}  \z{q_1} \zt{1}  \zt{1}  \zt{1}  \zb \zb \\
&\zb \zb \zt{1}  \zt{0}  \z{q_a} \zt{1}  \zt{1}  \zb \zb \\
\end{align*}
Nach der Verarbeitung steht also das Wort $\texttt{1011}_2=11_{10}$
auf dem Band.
\item
Die Verarbeitung des Wortes \texttt{1} ergibt
\begin{align*}
&\zb \zb \z{q_0} \zt{1}  \zb \zb \\
&\zb \zb \zt{1}  \z{q_0} \zb \zb \\
&\zb \zb \z{q_1} \zt{1}  \zb \zb \\
&\zb \zb \zt{0}  \z{q_a} \zb \zb \\
\end{align*}
Nach der Verarbeitung steht das Wort \texttt{0} auf dem Band.
\item
Die Verarbeitung des Wortes \texttt{0} ergibt
\begin{align*}
&\zb \zb     \z{q_0} \zt{0}  \zb \zb \\
&\zb \zb     \zt{0}  \z{q_0} \zb \zb \\
&\zb \zb     \z{q_1} \zt{0}  \zb \zb \\
&\zb \z{q_1} \zb     \zt{1}  \zb \zb \\
&\zb \zb     \z{q_r} \zt{1}  \zb \zb \\
\end{align*}
Die Null kann nicht dekrementiert werden.
\qedhere
\end{teilaufgaben}
\end{loesung}

\begin{bewertung}
Vollständiges Zustandsdiagramm ({\bf Z}) 1 Punkt,
nach rechts Fahren bis zum Ende ({\bf R}) 1 Punkt,
Dekrementoperation endet beim ersten \texttt{1} ({\bf E}) 1 Punkt,
Berechnungsgeschichte für a) ({\bf A}) 1 Punkt,
Berechnungsgeschichte für b) ({\bf B}) 1 Punkt,
Berechnungsgeschichte für c) ({\bf C}) 1 Punkt.
\end{bewertung}


